\begin{trapbox}

	\begin{description}[style=nextline, leftmargin=2.8em, labelsep=0.5em, itemsep=0.35em]
		\item[\textbf{\textcolor{CTrap}{易错点一（混淆角平分线）}}]
			错误认为在椭圆中切线平分角，在双曲线中法线平分角。
		\item[\textbf{\textcolor{CTrap}{正确理解（区分角色）：}}]
			椭圆在$P$处的法线平分$\angle F_1PF_2$；双曲线在$P$处的切线平分$\angle F_1PF_2$；抛物线在$P$处的切线平分$\angle FPH$。
		\item[\textbf{\textcolor{CTrap}{错因分析（记忆混淆）：}}]
			对三种圆锥曲线的光学性质记忆不准确，将法线与切线的角色混淆。
	\end{description}

	\medskip
	\begin{description}[style=nextline, leftmargin=2.8em, labelsep=0.5em, itemsep=0.35em]
		\item[\textbf{\textcolor{CTrap}{易错点二（忽略斜率不存在）}}]
			在验证角平分过程中，默认所有直线斜率都存在，未考虑焦半径或切线垂直于坐标轴的情况。
		\item[\textbf{\textcolor{CTrap}{正确理解（分类讨论）：}}]
			当出现斜率不存在时，应使用倾斜角或直接几何关系验证，性质在特殊点仍然成立。
		\item[\textbf{\textcolor{CTrap}{错因分析（缺乏分类）：}}]
			缺乏对垂直情况的分类讨论，导致推导不完整。
	\end{description}

	\medskip
	\begin{description}[style=nextline, leftmargin=2.8em, labelsep=0.5em, itemsep=0.35em]
		\item[\textbf{\textcolor{CTrap}{易错点三（抛物方向混淆）}}]
			在抛物线光学性质中，误认为法线平分$\angle FPH$，或误将$PH$取为指向焦点方向。
		\item[\textbf{\textcolor{CTrap}{正确理解（切线分角）：}}]
			抛物线在$P$处的切线平分$\angle FPH$，其中$PH$是过$P$平行于对称轴且指向开口方向的射线。
		\item[\textbf{\textcolor{CTrap}{错因分析（方向错误）：}}]
			对抛物线的反射特性理解不深，与椭圆性质混淆或未注意射线方向。
	\end{description}

\end{trapbox}
